\section{Final checklist}

This chapter introduced a new level of randomness. Earlier chapters studied
random observations and their probability laws. We now treat the entire sample
as a random object and study quantities computed from that sample.

The central chain of ideas is
\[
\text{probability model}
\longrightarrow
\text{random sample}
\longrightarrow
\text{statistic}
\longrightarrow
\text{sampling distribution}.
\]

\subsection*{Conceptual understanding}

You should be able to explain the following distinctions:

\begin{itemize}
    \item a random observation \(X\) versus an observed value \(x\);
    \item a random sample \((X_1,\ldots,X_n)\) versus an observed sample
    \((x_1,\ldots,x_n)\);
    \item a statistic \(T=g(X_1,\ldots,X_n)\) versus its observed value
    \(t=g(x_1,\ldots,x_n)\);
    \item the population distribution versus the law of the complete sample;
    \item the empirical distribution of one dataset versus the sampling
    distribution of a statistic;
    \item variability among individual observations versus variability among
    sample means or sample proportions.
\end{itemize}

\subsection*{Main definitions}

A random sample is a collection
\[
X_1,X_2,\ldots,X_n
\]
generated according to a specified sampling model.

A statistic is a function of the sample:
\[
T=g(X_1,\ldots,X_n).
\]

The sampling distribution of \(T\) is the probability law of that statistic over
all possible samples generated by the model.

\subsection*{Main examples}

For the sample mean,
\[
\overline X=\frac1n\sum_{i=1}^nX_i.
\]
If the observations have common expectation \(\mu\), then
\[
\E[\overline X]=\mu.
\]
If they are independent and have common variance \(\sigma^2\), then
\[
\Var(\overline X)=\frac{\sigma^2}{n}.
\]

For zero--one observations, the sample proportion is
\[
\widehat p=\frac1n\sum_{i=1}^nX_i.
\]
If \(P(X_i=1)=p\), then
\[
\E[\widehat p]=p,
\qquad
\Var(\widehat p)=\frac{p(1-p)}{n}.
\]

These formulas describe the centre and spread of sampling distributions. They do
not say that every observed statistic equals the corresponding population
quantity.

\subsection*{Typical tasks}

After studying this chapter, you should be able to:

\begin{itemize}
    \item identify the complete random sample in a simple model;
    \item list all possible samples when the sample space is small;
    \item assign probabilities to complete samples;
    \item compute a statistic for each possible sample;
    \item construct a sampling distribution by grouping samples that give the
    same statistic value;
    \item explain why a statistic is a random variable;
    \item distinguish a histogram of observations from a histogram of statistics
    obtained from repeated samples;
    \item explain why larger samples produce less variable sample means and
    sample proportions;
    \item describe how simulation approximates a sampling distribution.
\end{itemize}

\subsection*{Common mistakes}

Avoid the following mistakes:

\begin{itemize}
    \item treating the observed dataset as the only possible sample;
    \item confusing the empirical distribution with the sampling distribution;
    \item assuming that possible statistic values are equally likely;
    \item forgetting that probabilities of statistic values come from the
    probabilities of the samples producing them;
    \item interpreting \(\E[\overline X]=\mu\) as saying that every sample mean
    equals \(\mu\);
    \item claiming that the sample mean is normally distributed without an
    additional argument;
    \item treating a simulation as an exact proof of a probability law;
    \item forgetting that independence and a common distribution are modelling
    assumptions.
\end{itemize}

\subsection*{Questions for self-explanation}

A good test of understanding is whether you can answer these questions clearly:

\begin{enumerate}
    \item In what sense is a sample random before it is observed?
    \item What is the difference between \((X_1,\ldots,X_n)\) and
    \((x_1,\ldots,x_n)\)?
    \item Why does a statistic have a probability distribution?
    \item How is the probability of a statistic value obtained from the
    probabilities of complete samples?
    \item What is the difference between a population distribution, an empirical
    distribution, and a sampling distribution?
    \item Why do different samples from the same model produce different means?
    \item Why does the variance of \(\overline X\) decrease when \(n\) increases?
    \item Why is a sample proportion a special case of a sample mean?
    \item What does repeated sampling reveal that one observed dataset cannot?
\end{enumerate}

\begin{summarybox}
A sample is not only a collection of observed numbers. Before observation it is
a random object governed by a probability model. Every statistic computed from
that sample is therefore also random and has its own probability law. This new
level of randomness is the bridge from descriptive statistics to limit theorems
and statistical inference.
\end{summarybox}
